% =====================================================================
%  Chapter 22 — The Proof in Natural Language, 6
% =====================================================================

\chapter{The Budget}
\label{ch:budget2}

The chain buys multiplicity and it costs denominators. The budget is the
accounting of the two, and it is where the argument stops.

\section{The accounting}Recall the two parameters: the \emph{available growth} $c$, the multiplicative
factor by which the chain enlarges the space, and the \emph{budget} $q$, what
the segment of the argument actually reaches. The formalisation imposes:

\begin{center}
\begin{tabular}{@{}>{\RaggedRight\arraybackslash\hyphenpenalty=10000\exhyphenpenalty=10000}p{0.30\textwidth}>{\RaggedRight\arraybackslash\hyphenpenalty=10000\exhyphenpenalty=10000}p{0.62\textwidth}@{}}
\toprule
\textbf{Condition} & \textbf{Reading} \\
\midrule
$5 \le k$ & the weight is at least five \\
$1 \le N_s$ & the integral structure is non-trivial \\
$40\,(kN_s) \le c$ & the available growth must dominate forty times the
weight times the level cost \\
$\tfrac{9}{10}\,c \le q$ & the budget reached must be at least nine tenths
of the available growth \\
\bottomrule
\end{tabular}
\end{center}

The third condition is the one that says the multiplicativity must be strong
enough to pay for the integrality. The fourth says the argument must actually
get close to using all of what is available.

\section{The threshold}\begin{seamnote}{The threshold condition, in the source}
{\footnotesize
\begin{leanquote}
theorem threshold_arith_letter_case (k Ns c q bo : Real)
(hk5 : 5 <= k) (hN : 1 <= Ns) (hc : 40 * (k * Ns) <= c) (hq
      : 9 / 10 * c <= q)
    (hb : q - k < bo) : c / 2 <= 2 / 3 * bo - k * Ns
\end{leanquote}
}
\end{seamnote}

\begin{quote}
\textbf{Proposition Q (the threshold).} The argument succeeds if and only if
the available growth $c$ satisfies
\[
\frac{c}{2} \;\le\; \frac{2}{3}\cdot(\text{budget reached}) \;-\; kN_s.
\]
\end{quote}

Read the two sides.

\begin{description}[style=nextline,leftmargin=0pt,font=\sffamily\bfseries\small,
  itemsep=6pt]
\item[Left.] Half the available growth. This is the amount of growth the
argument needs in order to produce a square-root-quality remainder: the
remainder's size is the square root of the growth available, so half the
growth must be at least the square root of itself --- and the budget must
supply that.
\item[Right.] Two thirds of the budget the segment reached, less the
weight-and-level cost charged at each stage.
\end{description}

So the fraction of the budget the argument may use is two thirds.

\section{Why the third is lost}The two-thirds is not an arbitrary constant. It is the price of the atypical
path.

\begin{quote}
\textbf{Proposition R.} At each letter of a segment, the argument splits into
two cases: the \emph{typical} path, in which the letter's contribution is
tightly controlled by the per-letter budget, and the \emph{atypical} path, in
which it may exceed that budget by an amount of order the budget itself. The
atypical path must be paid for, and what survives of the budget after paying
for it is two thirds.
\end{quote}

The formal statement is \code{GoodA} in
\code{AtypicalPathSmallWeight.lean}: the atypical path satisfies the
relaxed bound
\[
\sum_c \bigl(\text{tangent parts}\bigr) \;\le\; \Xi + 2\cdot B(\mathrm{od}\ i),
\]
while the typical path \code{Typ} satisfies the tight bound without the slack.
And both carry the same requirement
\[
\frac{2}{3}\,B(\mathrm{od}\ i) \;\le\; \mathrm{plPart}\ \mu\ (\mathrm{lastIdx}\ r),
\]
which says: two thirds of the budget must be carried by the weight available at
the end of the segment.

\section{What would have to change}\begin{quote}
\textbf{Proposition S.} To obtain the whole budget rather than two thirds of
it, the argument would need
\[
\frac{c}{2} \;\le\; (\text{budget reached}) \;-\; kN_s,
\]
and the difference between the two right-hand sides is exactly one third of
the budget reached.
\end{quote}

That is the arithmetic of crossing the ceiling, and the remark's item 14 lists
three attempts to cross it:

\begin{center}
\begin{tabular}{@{}>{\RaggedRight\arraybackslash\hyphenpenalty=10000\exhyphenpenalty=10000}p{0.32\textwidth}>{\RaggedRight\arraybackslash\hyphenpenalty=10000\exhyphenpenalty=10000}p{0.62\textwidth}@{}}
\toprule
\textbf{Attempt} & \textbf{What it would buy} \\
\midrule
``get a main term'' & More of the available growth $c$ --- more lifts, or a
better per-lift multiplicativity. \\
``transfer constant to a function field'' & Move the constant to a setting
where the Weil machinery is native and the arithmetic factor exists. \\
``kill via twisting'' & Eliminate the atypical path outright, so that the
third is never spent. \\
``division'' & Something has to be divided; the development has a file named
for it. \\
\bottomrule
\end{tabular}
\end{center}

\section{The honest statement}\begin{quote}
\textbf{Proposition T (the obstruction, stated precisely).} The formalised
threshold condition holds with the factor $2/3$. Whether a variant with the
factor $1$ can be established --- by any of the three routes listed above, or
by a fourth not described in the remark --- is not settled by anything in the
development that has been audited, and the remark does not say.
\end{quote}

This is the honest formulation, and it is deliberately weaker than either
``the proof works'' or ``the proof fails''. What is established is the
computation of the ceiling; what is not established is its crossing. The
remark records the ceiling on 15 November and GRH on 16 November, and does not
say what happened in between.

\section{The reading against the earlier theorem}There is a strong reason to take the ceiling seriously rather than as an
artifact of the summary. The same author published, three months earlier,
arXiv:2608.13637, ``More than two thirds of the zeta zeros are simple and on
the critical line'' --- using the same method family, and formally verified in
Lean.

\begin{quote}
The most economical reading is that the budget computation above \emph{is} the
explanation of the two-thirds theorem: the method, costed, delivers a bound
covering two thirds of the zeros, which is exactly what the earlier paper
proves, and the three attempts are what was done about the remaining third.
\end{quote}

This is a hypothesis, and it is the most plausible one available; but it is a
hypothesis about the author's intent, and Chapter~\ref{ch:verdict} keeps it
distinct from what the artifact has been shown to do.

\begin{keytakeaways}
\item The budget condition is $c/2 \le \tfrac23\cdot(\text{budget reached})
- kN_s$: half the available growth must be covered by two thirds of the
budget, after the weight-and-level cost.
\item The missing third is the price of the atypical path, which may exceed
the per-letter budget by an amount of order the budget itself.
\item Crossing the ceiling requires one third more of the budget, which is
what the remark's three attempts --- main term, transfer of constants,
killing by twisting --- and its ``division'' are for.
\item The ceiling is computed and formalised; the crossing is not established
by anything audited, and the remark does not say.
\end{keytakeaways}
